\documentclass[11pt,a4paper]{report}
\usepackage[utf8]{inputenc}
\usepackage[T1]{fontenc}
\usepackage{cmbright}
\usepackage[brazil]{babel}
\usepackage[left=0cm,right=0cm,top=0cm,bottom=0cm]{geometry}
\usepackage{tikz,ulem,graphicx}
\usetikzlibrary{calc}
\pagestyle{empty}

\begin{document}
  \noindent\begin{tikzpicture}
  \fill (0,0) circle (0pt);
  \draw (1,-1) rectangle (20,-28.7);
  \draw (2,-2.7312) -- (2,-28.1229)
        (2.5771,-2.7312) -- (2.5771,-28.1229)
        (7,-1) -- (7,-2.7312)
        (7,-2.0387) -- (20,-2.0387)
        (1,-2.7312) -- (20,-2.7312)
        (1,-28.1229) -- (20,-28.1229);
  \node[right] at (7,-2.4139) {\LARGE\it Conteúdos};
    \node[inner sep=0pt] at (2.2885,-3.0198) {14};
    \node[right] at (2.5771,-3.0198) {\parbox{17.15cm}{\textbf{RECUPERAÇÃO PARALELA} $-$ Recuperação paralela do 3º período}};
    \draw (2,-3.3083) -- (20,-3.3083);
    \node[inner sep=0pt] at (2.2885,-3.5969) {14};
    \node[right] at (2.5771,-3.5969) {\parbox{17.15cm}{{\color{olive}\textbf{TÓPICOS DE GEOMETRIA} $-$ Relações métricas no triângulos retângulo, apresentação e conceitos.}}};
    \draw (2,-3.8854) -- (20,-3.8854);
    \node[inner sep=0pt] at (2.2885,-4.1740) {16};
    \node[right] at (2.5771,-4.1740) {\parbox{17.15cm}{{\color{olive}\textbf{TÓPICOS DE GEOMETRIA} $-$ Relações métricas no triângulos retângulo, dedução das relações.}}};
    \draw (2,-4.4625) -- (20,-4.4625);
    \node[inner sep=0pt] at (2.2885,-4.7510) {21};
    \node[right] at (2.5771,-4.7510) {\parbox{17.15cm}{\textbf{TÓPICOS DE GEOMETRIA} $-$ Relações métricas, exercícios propostos, questões 1, 2 e 3.}};
    \draw (2,-5.0396) -- (20,-5.0396);
    \node[inner sep=0pt] at (2.2885,-5.3281) {21};
    \node[right] at (2.5771,-5.3281) {\parbox{17.15cm}{{\color{olive}\textbf{TÓPICOS DE GEOMETRIA} $-$ Relações métricas, correção das questões 1, 2 e 3.}}};
    \draw (2,-5.6167) -- (20,-5.6167);
    \node[inner sep=0pt] at (2.2885,-5.9052) {28};
    \node[right] at (2.5771,-5.9052) {\parbox{17.15cm}{\textbf{TÓPICOS DE GEOMETRIA} $-$ Exercício proposto 4. Correção dos exercícios 3 e 4.}};
    \draw (2,-6.1937) -- (20,-6.1937);
    \node[inner sep=0pt] at (2.2885,-6.4823) {04};
    \node[right] at (2.5771,-6.4823) {\parbox{17.15cm}{\textbf{SISTEMA DE MEDIDAS} $-$ Medidas de comprimento e de área (aula remota)}};
    \draw (2,-6.7708) -- (20,-6.7708);
    \draw (1,-6.1937) -- (2,-6.1937);
    \node at (1.5,-4.4625) {\Large out};
    \draw (1,-6.7708) -- (2,-6.7708);
    \node at (1.5,-6.4823) {\Large nov};
  \node[right] at (1,-1.2885) {\large\bf\underline{SEDUC} / \underline{São Luis$-$MA}};
  \node[right] at (1,-1.8656) {\large\bf\underline{CE Poeta Cunha Santos}};
  \node[right] at (1,-2.4427) {\large\bf\underline{100 matutino}};
  \node[inner sep=0pt] at (13.5,-1.5771) {\bf\resizebox{12.5cm}{0.55cm}{Conteúdos de Letr Matemática do 4º Período / 2026}};
  \node[inner sep=0pt,above] at (10.5,-28.65) {Data: \underline{\,05\,}/\underline{\,11\,}/\underline{\,2026\,} \hspace{5mm} Aulas dadas: \underline{\,7\,} \hspace{5mm} Aulas previstas: \underline{\,7\,} \hspace{5mm} Professor(a): \underline{\includegraphics[width=0.2\linewidth]{../informados/.assinatura.png}}};
  \end{tikzpicture}
\end{document}
